\section{Introduction}
Frontier video-language models can answer a wide range of questions about complex scenes, but they frequently fail on \emph{physics checking}: detecting subtle inconsistencies such as discontinuous motion, identity swaps under occlusion, or uncaused changes in energy or momentum.
These failure modes matter in high-stakes settings (robotics, surveillance, simulation validation), where a system must not only describe what it sees but also flag implausible or impossible events.

Existing video understanding benchmarks provide broad coverage, but many tasks can be solved with weak temporal reasoning (e.g., static frames, language priors, or dataset artifacts), which reduces diagnostic value.
Conversely, physics-focused datasets often emphasize \emph{prediction} (``what happens next?'') rather than \emph{verification} (``did an impossible event occur?'') or they focus on a narrow set of physical principles.
Recent work on evaluating video \emph{generation} models as world models highlights the importance of physics adherence, but typically evaluates \emph{generated} videos and uses human preference labels \cite{li2025worldmodelbench}.
Our focus is complementary: we evaluate \emph{video reasoning} models on whether they can \emph{detect} minimal physical inconsistencies in observed clips.

\paragraph{Benchmark design principles.}
MIRAGE-V is built around three principles:
(1) \textbf{Minimal-change pairs} --- each datapoint is part of an $A/B$ pair that shares context but differs by a single intervention, enabling causal attribution and the \emph{pair accuracy} metric;
(2) \textbf{Taxonomy-driven failures} --- interventions are mapped to a small set of interpretable modules (taxonomy categories) covering temporal, occlusion, latent-parameter, invariant, and causal phenomena;
(3) \textbf{Sampling-robust evaluation} --- each example provides a critical window when decisive evidence occurs, supporting both sparse and dense sampling protocols.

\paragraph{Contributions.}
We contribute:
(i) \textbf{MIRAGE-V}, a benchmark suite of minimal-change video pairs with taxonomy labels and critical-window annotations;
(ii) an \textbf{open-source generator} that produces scalable synthetic clips and optional real-video edits aligned to the taxonomy;
(iii) a \textbf{unified evaluation harness} for open-weight and API-based models with standardized prompting, metrics, and shortcut audits;
(iv) an empirical baseline on three representative models demonstrating that physics checking remains challenging, especially under sparse frame sampling.
